%!TEX TS-program = lualatex
%!TEX encoding = UTF-8 Unicode

\documentclass[t]{beamer}

%%%% HANDOUTS For online Uncomment the following four lines for handout
%\documentclass[t,handout]{beamer}  %Use this for handouts.
%\usepackage{handoutWithNotes}
%\pgfpagesuselayout{3 on 1 with notes}[letterpaper,border shrink=5mm]

%\includeonlylecture{student}

%%% Including only some slides for students.
%%% Uncomment the following line. For the slides,
%%% use the labels shown below the command.

%% For students, use \lecture{student}{student}
%% For mine, use \lecture{instructor}{instructor}

% FONTS
\usepackage{fontspec}
\def\mainfont{Linux Biolinum O}
\setmainfont[Ligatures={Common,TeX}, Contextuals={NoAlternate}, BoldFont={* Bold}, ItalicFont={* Italic}, Numbers={OldStyle}]{\mainfont}
\setsansfont[Ligatures={Common,TeX}, Scale=MatchLowercase, Numbers=OldStyle, BoldFont={* Bold}, ItalicFont={* Italic},]{Linux Biolinum O} 
\usepackage{microtype}

\usepackage{graphicx}
	\graphicspath{{/Users/mtaylor/pictures/teach/163/lecture/}
	{/Users/mtaylor/pictures/teach/common/}} % set of paths to search for images

\usepackage{multicol}
%\usepackage{booktabs}
%\usepackage{textcomp}

\usepackage{tikz}
	\tikzstyle{every picture}+=[remember picture,overlay]
%	\usetikzlibrary{arrows}

\newcommand{\blurred}[3]{
\draw [black!#1] node at (#2,#2) {#3};
\draw [black!#1] node at (0,#2) {#3};
\draw [black!#1] node at (#2,0) {#3};
\draw [black!#1] node at (-#2,#2) {#3};
\draw [black!#1] node at (-#2,0) {#3};
\draw [black!#1] node at (0,-#2) {#3};
\draw [black!#1] node at (#2,-#2) {#3};
\draw [black!#1] node at (-#2,-#2) {#3};
}

\newcommand{\blurry}[1]{
\blurred{5}{0.04}{#1}
\blurred{13}{0.017}{#1}
\blurred{10}{0.012}{#1}
\blurred{15}{0.007}{#1}
\blurred{20}{0.003}{#1}
\blurred{10}{0.0015}{#1}
\draw [black!40] node at (0,0) {#1};
}

\mode<presentation>
{
  \usetheme{Lecture}
  \setbeamercovered{invisible}
}


\begin{document}
%
\lecture{student}{student}
%
\begin{frame}{Our goal for this lecture is to }
	
	\hangpara explain four of the five processes that cause evolution in populations:
	
	\hangpara \quad \highlight{mutation,}

	\hangpara \quad \highlight{gene flow,}
	 	 
	\hangpara \quad \highlight{non-random mating,} and
	
	\hangpara \quad \highlight{genetic drift.} 

\end{frame}
% 
\begin{frame}[t]{The Hardy-Weinberg equations are a \highlight{null model} for population evolution.}
	
	\vspace{-\baselineskip}
	
	\hangpara The H-W equations makes five assumptions.
	
	\hangpara \quad Mutations do not occur.
		
	\hangpara \quad Gene flow does not occur.

	\hangpara \quad Mating is random.
	
	\hangpara \quad Genetic drift does not occur.
	
	\hangpara \quad Natural selection does not occur.
	
	\pause \hangpara \highlight{If \emph{at least one} of these occurs in a population, then the population \emph{will} evolve.}
	
\end{frame}
%
{
\usebackgroundtemplate{\includegraphics[width=\paperwidth]{mutation_types} }
\begin{frame}[t]{A \highlight{mutation} is any change to the genome of an organism.}

	\vfilll

	\hfill\rotatebox{90}{\tiny \textcopyright Roberts \& Company}
\end{frame}
}
%
\begin{frame}[t]{Mutations can create new alleles.}

	\hangpara Mutations occur at \highlight{random.}\pause
	
	\hangpara Most mutations are \highlight{neutral mutations.}\pause
	
	\hangpara Mutations that occur in a gene can create new alleles.\pause
	
	\hangpara Some new alleles can be detrimental or beneficial.\pause

	\hangpara More alleles increase \highlight{heterozygosity} in a population.\pause
	
	\hangpara Higher heterozygosity increases genetic variation for natural selection to work with.
	
\end{frame}
%
\begin{frame}[t]{\highlight{Gene flow} is the movement of individuals among populations.}
	\alt<handout>{}{\includegraphics<1>[width=\textwidth]{gene_flow_beetles_start}}
	\alt<handout>{}{\includegraphics<2>[width=\textwidth]{gene_flow_beetles_middle}}
	\includegraphics<3>[width=\textwidth]{gene_flow_beetles_end}

	\onslide<3>When individuals move, their alleles go with them.
	
	\vfilll
	
	\hfill \tiny Tsaneda, Wikimedia, \ccby{3}
\end{frame}
%
{
\usebackgroundtemplate{\includegraphics[width=\paperwidth]{mating_mice} }
\begin{frame}[b]{\highlight{Non-random mating} occurs when individuals choose mates based on phenotype.}
	\hfill \tiny \textcopyright Pearson Education, Inc.
\end{frame}
}
%
{
\usebackgroundtemplate{\includegraphics[width=\paperwidth]{charles_ancestry_inbred} }
\begin{frame}[b]{Non-random mating tends to reduce heterozygosity.}
\end{frame}
}
%
{
\usebackgroundtemplate{\includegraphics[width=\paperwidth]{shirt_sniffing}}
\begin{frame}[b]{Non-random mating can \emph{increase} heterozygosity for some genes.}
	\tiny \textcolor{white}{\copyright\,Associated Press}
\end{frame}
}
%
\begin{frame}[t]{\highlight{Genetic drift} is \emph{random} change of allele frequencies.}
	\vspace*{-\baselineskip}
	\begin{center}
		\includegraphics[width=0.7\textwidth]{genetic_drift}
	\end{center}

	\vspace*{-\baselineskip}
	
	\hangpara Genetic drift is strongest in small populations.

\end{frame}
%
\begin{frame}{Genetic drift \emph{always} causes a \highlight{loss of heterozygosity.}}
	\includegraphics[width=\textwidth]{loss_of_heterozygosity}
\end{frame}
%
\begin{frame}{Two types of drift cause \highlight{rapid} loss of heterozygosity.}

	\begin{multicols}{2}
		\includegraphics[width=0.5\textwidth]{loss_of_heterozygosity}
	\columnbreak
		\hangpara \highlight{Bottleneck effect} and\\ 
		\highlight{Founder effect.}
	\end{multicols}
	
\end{frame}
%
{
\usebackgroundtemplate{\includegraphics[width=\paperwidth]{bottleneck} }
\begin{frame}[b]{\highlight{Bottlenecks} occur when population size decreases rapidly.}

	\hfill \tiny \textcopyright Pearson Education, Inc.
\end{frame}
}
%
\lecture{instructor}{instructor}
%
{
\usebackgroundtemplate{\includegraphics[width=\paperwidth]{pacific_ocean} }
\begin{frame}[b]{}
\end{frame}
}
%
{
\usebackgroundtemplate{\includegraphics[width=\paperwidth]{pingelap_atoll} }
\begin{frame}[b]{}
	\hfill \tiny \textcolor{white}{Google Earth}
\end{frame}
}
%
{
\usebackgroundtemplate{\includegraphics[width=\paperwidth]{tropical_island_normal} }
\begin{frame}[b]{}
	\tiny Benedict Adam, Flickr, \ccby{2}.
\end{frame}
}
%
{
\usebackgroundtemplate{\includegraphics[width=\paperwidth]{tropical_island_blurred} }
\begin{frame}[b]{}
	\begin{tikzpicture}
		\blurry{\tiny\textcolor{gray}{Benedict Adam, Flickr, \ccby{2}}}
	\end{tikzpicture}
\end{frame}
}
%
\lecture{student}{student}
%
{
\usebackgroundtemplate{\includegraphics[width=\paperwidth]{pingelap_atoll} }
\begin{frame}[b]{\textcolor{white}{1 in 17 people on Pingelap have achromatopsia.}}
	\hfill \tiny \textcolor{white}{Google Earth}
\end{frame}
}
%
{
\usebackgroundtemplate{\includegraphics[width=\paperwidth]{founder_effect_silver_eye} }
\begin{frame}[b,plain]{}
	\tiny \textcopyright Pearson Education \hfill Toby Hudson, Wikimedia, \ccbysa{3}
\end{frame}
}
%
\lecture{instructor}{instructor}
%
{
\usebackgroundtemplate{\includegraphics[width=\paperwidth]{hms_bounty} }
\begin{frame}[b]{}
	\hfill \textcolor{white}{\tiny Yasmina, Wikimedia, \ccby{3}}
\end{frame}
}
%
{
\usebackgroundtemplate{\includegraphics[width=\paperwidth]{tahiti} }
\begin{frame}[b]{}
	\hfill \textcolor{white}{\tiny Moises Gonzalez, Flickr, ccbync{2}}
\end{frame}
}
%
{
\usebackgroundtemplate{\includegraphics[width=\paperwidth]{bounty_mutiny} }
\begin{frame}[b]{}
	\hfill \textcolor{white}{\tiny National Maritime Museum, Wikimedia, public domain.}
\end{frame}
}
%
{
\usebackgroundtemplate{\includegraphics[width=\paperwidth]{mutineer_route} }
\begin{frame}[b]{}
	\tiny Benton et al. 2015. Investigative Genetics 6: 12.
\end{frame}
}
%
\end{document}
